% ---------------------------------------------------------------------------
\section{Reproduction of the Paper's Models (RQ1)}
\label{sec:reproduction}

Before comparing new models, we rebuilt the three models of
Severson~\textit{et al.}~\cite{severson2019} on our cleaned data. If our
pipeline reproduces their Table~1, any later difference between models can be
trusted to come from the models, not from the data handling.

\subsection{Features}
All 20 candidate features of the paper's Supplementary Table~1 are computed
by \texttt{src/features.py} from cycles $\le 100$ only (Table~\ref{tab:featsets}).
The three models use the subsets marked in that table:
\begin{itemize}
  \item \textbf{Variance} (1 feature): $\log_{10}\mathrm{Var}(\dq)$.
  \item \textbf{Discharge} (6): $\log_{10}$ of $|\min|$, variance, $|$skewness$|$
  and $|$kurtosis$|$ of $\dq$; discharge capacity at cycle 2; maximum
  capacity minus capacity at cycle 2.
  \item \textbf{Full} (9): $\log_{10}|\min|$ and $\log_{10}\mathrm{Var}$ of $\dq$;
  slope and intercept of a straight-line fit to capacity over cycles 2--100;
  capacity at cycle 2; mean charge time over cycles 2--6; temperature
  integral $\int T\,dt$ over cycles 2--100; minimum IR over cycles 2--100;
  IR(100)$-$IR(2).
\end{itemize}

\begin{table}[t]
\centering
\caption{Implementation choices where the paper's Supplementary Note~1 is
ambiguous.}
\label{tab:featsets}
\setlength{\tabcolsep}{3pt}
\begin{tabular}{@{}p{0.30\columnwidth}p{0.64\columnwidth}@{}}
\toprule
Item & Choice \\
\midrule
Variance & sample variance ($p-1$), as in Note~1 \\
Skewness, kurtosis & standard moment ratios $m_3/m_2^{1.5}$, $m_4/m_2^2$; the
Note~1 skewness denominator is read as a typo \\
Capacity features & \texttt{QD\_clean} (glitch cycles removed, rule B3) \\
Single-cycle values & nearest cycle with a finite value within $\pm3$ \\
Temperature integral & sum over cycles of the trapezoidal $\int T\,dt$ of the
raw traces, leaving out sample gaps $>10$\,min (rule B10) \\
Missing values & training-set median (only \texttt{b1c18}'s IR, rule B9) \\
\bottomrule
\end{tabular}
\end{table}

\subsection{Protocol}
As in the paper, the target is $\log_{10}$(cycle life), and every model is fitted
and tuned on the 41 training cells only. The variance model is an ordinary
least-squares line (the paper applies no model selection to it). The other
two standardise the features and fit an elastic net, with $\alpha$ and the
$\ell_1$ ratio chosen by 4-fold cross-validation repeated 10 times (the paper
uses 4-fold cross-validation with Monte Carlo sampling). The full model leaves
out the four cells whose temperature probe lost contact (rule B6), as the
paper does. The primary and secondary test sets are evaluated once. RMSE and
mean percentage error are computed on cycles, not on the log scale. The
constant baseline predicts the training mean of $\log_{10}$(life) for every
cell, i.e.\ the geometric mean; this reproduces the paper's Supplementary
Table~3 exactly, whereas the arithmetic mean does not.

\subsection{Results}
Table~\ref{tab:repro} compares our errors with the paper's. The
\textbf{variance model matches to within one cycle} on every split, which
confirms that our $\dq$, cell labels and splits are identical to the paper's.
The discharge model is within 1--9\,\% of the paper's RMSE. The full model is
close on the test sets (primary 125 vs.\ 118 cycles, secondary 220 vs.\ 214)
and less close on training (67 vs.\ 51). Its higher primary RMSE without the
anomalous cell (122 vs.\ 100) comes mostly from a single cell: \texttt{b1c2},
the longest-lived cell in the dataset (2237 cycles), is predicted at 1594
cycles and accounts for 67\,\% of the squared error. Without it the RMSE
would be 71 cycles. The ranking of the three models is the same as in the paper on every
split, including the paper's counter-intuitive result that the full model is
\emph{worse} than the discharge model on the secondary test set. The paper
attributes this to the secondary batch's extra year of calendar ageing,
which lowers its initial capacity and biases the capacity features.

\begin{table*}[t]
\centering
\caption{Reproduction of the paper's Table~1 and the constant baseline of
its Supplementary Table~3. Each entry is RMSE (cycles) / mean percentage
error; paper values in parentheses. ``Primary excl.'' leaves out the anomalous
cell \texttt{b2c1}. The policy-mean baseline is ours.}
\label{tab:repro}
\begin{tabular}{@{}lcccc@{}}
\toprule
Model & Train & Primary test & Primary excl. & Secondary test \\
\midrule
Constant      & 327 / 29.6\,\% (327 / 29.6) & 401 / 35.0\,\% (401 / 34.9) & 399 / 28.2\,\% (399 / 28.2) & 511 / 36.1\,\% (511 / 36.1) \\
Policy mean   & 46 / 1.6\,\%  & 173 / 27.3\,\% & 159 / 20.3\,\% & 509 / 36.0\,\% \\
Variance      & 104 / 14.1\,\% (103 / 14.1) & 138 / 14.7\,\% (138 / 14.7) & 138 / 13.2\,\% (138 / 13.2) & 196 / 11.4\,\% (196 / 11.4) \\
Discharge     & 69 / 9.2\,\% (76 / 9.8)     & 90 / 12.6\,\% (91 / 13.0)   & 85 / 9.4\,\% (86 / 10.1)    & 167 / 9.9\,\% (173 / 8.6) \\
Full          & 67 / 6.8\,\% (51 / 5.6)     & 125 / 11.0\,\% (118 / 14.1) & 122 / 8.0\,\% (100 / 7.5)   & 220 / 11.2\,\% (214 / 10.7) \\
\bottomrule
\end{tabular}
\end{table*}

\begin{figure*}[t]
  \centering
  \includegraphics[width=\textwidth]{fig_p2_parity}
  \caption{Predicted vs.\ observed cycle life for the three reproduced models
  (log axes; cf.\ the paper's Fig.~3). Dashed line: perfect prediction.
  RMSE in parentheses excludes the anomalous cell \texttt{b2c1} (circled). In
  all three models, the longest-lived secondary-test cells are under-predicted.}
  \label{fig:parity}
\end{figure*}

\textbf{Coefficients.} The paper prints the full model's coefficients on
standardised features (its Supplementary Fig.~5). Ours have the same sign for
all nine features and similar magnitudes (Table~\ref{tab:coef}). Like the
paper, our elastic net gives the fade-curve intercept a weight of essentially
zero. It also zeroes the fade-curve slope, to which the paper gives a small
weight.

\begin{table}[t]
\centering
\caption{Full-model coefficients on standardised features (target
$\log_{10}$ life): ours vs.\ the paper's Supplementary Fig.~5.}
\label{tab:coef}
\begin{tabular}{@{}lrr@{}}
\toprule
Feature & Ours & Paper \\
\midrule
$\log_{10}\mathrm{Var}(\dq)$ & $-0.049$ & $-0.071$ \\
$\log_{10}|\min\dq|$ & $-0.040$ & $-0.029$ \\
Charge time, cycles 2--6 & $+0.033$ & $+0.039$ \\
Temperature integral & $-0.029$ & $-0.036$ \\
Capacity, cycle 2 & $+0.025$ & $+0.022$ \\
Minimum IR & $-0.021$ & $-0.007$ \\
IR(100)$-$IR(2) & $-0.000$ & $-0.013$ \\
Fade slope, cycles 2--100 & $0.000$ & $-0.013$ \\
Fade intercept, cycles 2--100 & $0.000$ & $0.000$ \\
\bottomrule
\end{tabular}
\end{table}

\textbf{Stability.} Repeating the cross-validation with ten different random
seeds changes the test RMSE by only a few cycles (discharge: primary 89--91,
secondary 167--170; full: primary 117--128, secondary 218--222). The
remaining differences from the paper are of the same size as this spread.
They are expected given the paper's different software (MATLAB), its own
random sampling, and its spline fit of the discharge curves.

\textbf{Two data problems found while reproducing.} The first full-model fit
predicted 0 cycles for \texttt{b1c18} (primary RMSE 162). The cause was a
fault in this cell's records for cycles 2--53: a corrupt time stamp that
jumps by $\sim$4000\,min within every record, inflating the temperature
integral ($z = 420$), and an IR logger stuck at one value. Rules B9 and B10
(Table~\ref{tab:cleaning}) fix both. This shows why a reproduction is worth
doing before any new model: the paper's numbers served as a check that
exposed a data fault that the earlier cleaning had missed.

\textbf{Baselines.} The policy-mean baseline (our addition) does better than the
constant on the primary set (173 vs.\ 401 cycles), since the primary cells
share charging policies with training. But it has nothing to offer on the
secondary set, whose policies never appear in training, and it is still
clearly worse than the single-feature variance model. Both baselines confirm
that the $\dq$ features carry the predictive information.

\textbf{Answer to RQ1.} Yes: the paper's results are reproduced, exactly for
the variance model and to within the cross-validation spread for the others,
with the same model ranking. The pipeline is therefore a sound base for the
model comparison.
